\chapter{The Developer Tooling Ecosystem}

\section{The Super-Swiss-Army Knife}

\begin{kidbox}{The Complete Adventure Backpack}
When you go camping in the woods, you don't just take a pair of shoes.
You bring a tent, a compass, a warm sleeping bag, a flashlight, a water bottle, and a campfire kit!
Shae does not make you download dozens of complicated mystery tools from thirty different websites.
Every single tool you need to write, test, format, check, benchmark, document, and package your software is built directly into the \texttt{shae} command!
\end{kidbox}

\section{The Complete Tooling Suite}

\begin{longtable}{ll}
\toprule
\textbf{Subcommand} & \textbf{Purpose and Features} \\
\midrule
\endhead
\texttt{shae run <file.shae>} & Executes source files with \texttt{--engine=vm} or \texttt{--engine=ast} \\
\texttt{shae ast <file.shae>} & Pretty-prints the Abstract Syntax Tree for debugging \\
\texttt{shae disasm <file.shae>} & Disassembles bytecode into human-readable assembly \\
\texttt{shae fmt <file.shae>} & Formats code automatically to canonical Shae style guidelines \\
\texttt{shae fmt --check} & Verifies code formatting in continuous integration pipelines \\
\texttt{shae lint <file.shae>} & Detects unused bindings, dead code, and variable shadowing \\
\texttt{shae doc [path]} & Generates JSON API specifications and documentation \\
\texttt{shae doc --serve} & Spawns a local HTTP server serving documentation \\
\texttt{shae bench [filter]} & Runs microbenchmarks with high-resolution latency tracking \\
\texttt{shae llvm <file.shae>} & Compiles to native binary using LLVM SSA and Clang \\
\texttt{shae link <file.shaec>} & Links bytecode or object chunks into native ELF/PE binaries \\
\texttt{shae pkg init <name>} & Initializes a new package with manifest file (\texttt{shae.toml}) \\
\texttt{shae pkg build} & Compiles package source modules and dependencies \\
\texttt{shae pkg publish} & Publishes package to the central package registry \\
\bottomrule
\end{longtable}
